\section{数值结果}

图~\ref{fig:L_dependence} 展示了准 TMDWF 的模随 Wilson 线长度 $\ell$ 的依赖。
如图所示，在 $\{P^z, b_{\perp}, t\}=\{6\pi/L, 3a, 6a\}$ 情形下，准 TMDWF $\phi_\ell(0,b_\perp,P^z,\ell)$ 与 Wilson 圈 $Z_E$ 的平方根都随长度 $\ell$ 指数衰减，但减除后的准 TMDWF 在 $\ell\ge 0.4$ fm 时与长度无关。更大 $P^z$、$b_{\perp}$、$t$ 的其他情形见补充材料~\cite{supplemental}。基于这一观察，下面的计算中我们对所有情形取 $\ell=7a=0.686$ fm 作为渐近结果。

\begin{figure}[!th]
\begin{center}
\includegraphics[width=0.45\textwidth]{images/03_ff_P3b3.pdf}
\caption{比值 $C_3(b_\perp,P^z,t_{\rm sep},t)/C_2(0, P^z,0,t_{\rm sep})$（数据点）作为 $t_{\rm sep}$ 与 $t$ 的函数，在 $\{P^z, b_{\perp}\}=\{6\pi/L, 3a\}$ 下收敛于 $t,t_{\rm sep}\rightarrow \infty$ 处的基态贡献（灰带）。正如此图所示，我们的数据总体上与预期的拟合函数（彩色带）相符。}\label{fig:joint_fit}
\end{center}
\end{figure}

我们对相同 $P^z$ 与 $b_{\perp}$ 下的形状因子和准 TMDWF 按 Eqs.~\eqref{eq:3pt} 与 \eqref{eq:2pt} 的参数化做了联合拟合。$\{P^z, b_{\perp}\}=\{6\pi/L, 3a\}$ 情形下不同 $t_{\rm sep}$、$t$ 的比值 $C_3(b_\perp,P^z,t_{\rm sep},t)/C_2(0, P^z,0,t_{\rm sep})$ 示于图~\ref{fig:joint_fit}，图中同时给出基态贡献（灰带）与有限 $t_{2}$、$t$ 处的拟合结果（彩色带）。在本计算中，激发态贡献被拟合很好地描述，$\chi^2/{\rm d.o.f.}=0.6$。联合拟合的细节及更多拟合质量检验见补充材料~\cite{supplemental}，拟合质量类似。

\begin{figure}[!th]
\includegraphics[width=0.45\textwidth]{images/04_sf_perturb_sqrt2.pdf}
\caption{ 按式~\eqref{eq:S_ratio_z=0} 取 $b_{\perp,0}=a$ 时内禀软因子随 $b_{\perp}$ 的变化。不同π介子动量 $P^z$ 的结果彼此一致。虚线为一圈计算结果，见式~(\ref{eq:S_ratio_1loop})，强耦合常数取 $\alpha_s(1/b_{\perp})$。{阴影带对应 $\alpha_s$ 的标度不确定度：$\mu\in [1/\sqrt{2},\sqrt{2}]\times 1/b_\perp$。来自算符混合的系统不确定度已被计入。}}\label{fig:quasi-soft-factor}
\end{figure}

所得软因子随 $b_{\perp}$ 的变化绘于图~\ref{fig:quasi-soft-factor}，对应 $\gamma$= 2.17、3.06 与 3.98，即 $P^z=\{4,6,8\}\pi/L=\{1.05, 1.58, 2.11\}$ GeV。
%
如图~\ref{fig:quasi-soft-factor} 所示，不同大 $\gamma$ 下的结果彼此一致，
表明渐近极限在误差范围内是稳定的。
%
我们还将从格点提取的内禀软函数与式~(\ref{eq:S_ratio_1loop}) 的一圈结果比较，$\alpha_s(\mu=1/b_\perp)$ 由 $\alpha_s(\mu=2\;{\rm GeV})\approx 0.3$ 演化而来。{阴影带对应 $\alpha_s$ 的标度不确定度：$\mu\in [1/\sqrt{2},\sqrt{2}]\times 1/b_\perp$。}
%
注意前者的 $b_\perp$ 依赖完全来自格点模拟，而后者来自微扰论。
{为便于比较，我们还在补充材料~\cite{supplemental}中列表给出了软函数的结果。}


\begin{figure}[!th]
\begin{center}
\includegraphics[width=0.45\textwidth]{images/05a_phi2mod.pdf}
\includegraphics[width=0.45\textwidth]{images/05b_CS_kernel.pdf}
\caption{ 准 TMDWF（上图）与提取的 Collins-Soper 核（下图）随 $b_{\perp}$ 的变化。准 TMDWF 可见的 $P^z$ 依赖主要可以由 Collins-Soper 核来理解，因为我们用树级匹配得到的核与三圈微扰结果（在小 $b_{\perp}$ 处）以及π介子准 TMDPDF 的非微扰结果在标度 $\alpha_s(1/b_{\perp})$ 下一致。{基于 quenched 格点计算、标记为 ``Hermite'' 与 ``Bernstein''~\cite{Shanahan:2020zxr} 的结果也一并给出以供比较。下图的误差包含统计误差以及来自非零虚部与算符混合效应的系统误差。}}\label{fig:quasi-TMDWF}
\end{center}
\end{figure}

如图~\ref{fig:quasi-TMDWF} 上图所示，用 $\phi_\ell(0,0,P^z,0)$ 归一化的准 TMDWF $|\phi_\ell(0,b_{\perp},P^z,\ell)|$ 呈现出清晰的 $P^z$ 依赖。按式~(\ref{eq:CS_kernel_z=0})，这一依赖与 CS 核相关联，只带有可能的 LaMET 匹配效应与 $1/\gamma^2$ 量级的幂次修正。因此我们用式~(\ref{eq:CS_kernel_z=0}) 在树级近似下提取该核，并将图~\ref{fig:quasi-TMDWF} 下图的结果与文献~\cite{Shanahan:2020zxr} 以及取 $\alpha_s(\mu=1/b_\perp)$ 的最高三圈微扰结果比较。{我们通过将算符混合效应的贡献与准 TMDWF 非零虚部的贡献【后者应被{对高阶效应的适当处理所抵消}】平方相加（in quadrature）来估计系统不确定度。}细节见补充材料~\cite{supplemental}，特别是{Sec.~\ref{subsec:mixing} 与 \ref{subsec:imagpart}}。我们的结果与文献~\cite{Shanahan:2020zxr}一致。
%, but the results based on $P^z_1/P^z_2=3/2$ and $P^z_1/P^z_2=4/2$ have some differences at certain $b_{\perp}$ which might come from insufficient statistics at $P^z$=2.11 GeV, or potential  systematic uncertainties.
